\documentclass[11pt,a4paper]{article}

\usepackage[utf8]{inputenc}
\usepackage[T1]{fontenc}
\usepackage{amsmath,amssymb}
\usepackage{geometry}
\usepackage{booktabs}
\usepackage{longtable}
\usepackage{array}
\usepackage{hyperref}
\usepackage{enumitem}
\usepackage{caption}

\geometry{margin=2.5cm}
\hypersetup{colorlinks=true, linkcolor=black, citecolor=black, urlcolor=black}

\newcolumntype{L}[1]{>{\raggedright\arraybackslash}p{#1}}

\title{Fiche technique --- Contribution du Membre 5\\
\large Monte-Carlo, implémentation et expérimentation}
\author{Projet : distribution équitable de l'eau sous demande incertaine}
\date{}

\begin{document}
\maketitle

\begin{quote}
\small
Ce document couvre exclusivement la contribution du Membre 5 (M5). Il ne
reprend ni ne redérive les parties théoriques complètes de M1, M2, M3 ou M4 ;
il n'en présente que ce qui est strictement nécessaire à la compréhension de
l'implémentation de M5. Toute information absente des fichiers transmis est
signalée explicitement par la mention \emph{« Information non disponible dans
les fichiers fournis »} plutôt que supposée ou inventée. Ce document n'est pas
le rapport final du projet.
\end{quote}

\tableofcontents

% =====================================================================
\section{Présentation de la contribution de M5}
% =====================================================================

M5 est responsable du volet « Monte-Carlo, implémentation et
expérimentation » du projet. Concrètement, cette responsabilité couvre :

\begin{itemize}
  \item l'implémentation du simulateur Monte-Carlo de scénarios de demande
  (\texttt{src/probability/monte\_carlo.py}) ;
  \item l'implémentation de la fonction de coût pénalisée et de son gradient
  (\texttt{src/optimization/objective.py}), à partir des dérivations
  mathématiques fournies par M4 ;
  \item l'implémentation du solveur de descente de gradient projetée
  (\texttt{src/optimization/gradient\_descent.py}), également à partir des
  dérivations de M4 ;
  \item l'orchestration de la résolution du problème d'optimisation sur un
  ensemble de scénarios (\texttt{src/simulation/run\_scenarios.py}) ;
  \item la réalisation de plusieurs expériences visant à confronter les
  résultats théoriques établis par M2 et M4 au comportement observé du code.
\end{itemize}

M5 implémente ; M5 ne dérive pas la théorie. Les formules utilisées dans
\texttt{objective.py} et \texttt{gradient\_descent.py} proviennent d'un
document rédigé par M4 (\emph{Aïchatou Traoré}),
\texttt{docs/section\_optimisation\_membre4.pdf}, cité nommément dans les deux
fichiers. Le rôle de M5 est de traduire ces dérivations en code correct,
vérifié, et de produire les expériences qui les confrontent au comportement
réel du système.

Un point d'audit important est signalé dès cette section et repris en détail
en section~15 : le fichier \texttt{exp6\_at\_risk\_districts.py}, bien que
présent dans le périmètre nominal de M5 selon la numérotation des six
expériences du projet, s'auto-attribue explicitement à M3 dans son propre
en-tête et n'utilise aucun des quatre modules de M5. Il est documenté ici tel
qu'il a été transmis, avec cette réserve clairement posée.

% =====================================================================
\section{Position de M5 dans l'architecture globale}
% =====================================================================

L'architecture du projet est organisée comme suit :

\begin{verbatim}
data/
src/graph/          (M1, M2)
src/probability/    (M3 ; monte_carlo.py appartient à M5)
src/optimization/   (M5, à partir des dérivations de M4)
src/simulation/      (M5)
src/evaluation/      (hors périmètre de M5)
experiments/         (scripts 3, 4, 5 attribués à M5 ; script 6 voir section 15)
results/
tests/
report/
\end{verbatim}

Le tableau~\ref{tab:pipeline} récapitule les étapes du pipeline telles que
mentionnées dans les en-têtes des fichiers fournis, avec le responsable
associé et le fichier correspondant.

\begin{longtable}{L{1.5cm}L{4.5cm}L{2.5cm}L{5cm}}
\caption{Étapes du pipeline mentionnées dans les fichiers fournis.}
\label{tab:pipeline}\\
\toprule
Étape & Description & Responsable & Fichier \\
\midrule
\endhead
3 & Constante de Lipschitz / conditionnement (référencée, non fournie) & M2 & \texttt{src/graph/graph\_analysis.py} (non fourni) \\
4 & Modèle de demande $D_i \sim \mathcal{N}(\mu_i,\sigma_i^2)$ & M3 & \texttt{demand\_model.py} \\
5 & Génération de scénarios de demande & M5 (avec M3) & \texttt{monte\_carlo.py} \\
6 & Formulation du problème et convexité & M4 & \texttt{objective.py} (implémentation M5) \\
7 & Pénalisation de la contrainte d'égalité & M4 & \texttt{objective.py} (implémentation M5) \\
8 & Gradient $\nabla J(q)$ & M4 & \texttt{objective.py} (implémentation M5) \\
9 & Descente de gradient projetée & M5 (à partir de M4) & \texttt{gradient\_descent.py} \\
10 & Résolution sur un ensemble de scénarios & M5 & \texttt{run\_scenarios.py} \\
\bottomrule
\end{longtable}

% =====================================================================
\section{Objectifs de l'implémentation}
% =====================================================================

D'après les en-têtes des fichiers fournis, l'implémentation de M5 poursuit
trois objectifs explicites :

\begin{enumerate}
  \item Transformer le modèle probabiliste de M3 en scénarios de demande
  exploitables, avec une mesure explicite de la précision statistique des
  quantités estimées (voir le docstring de \texttt{monte\_carlo.py}, qui
  insiste sur ce point : \emph{« Ce qui est noté, ce n'est pas de savoir
  tirer des nombres au hasard, c'est de savoir avec quelle précision on
  connaît les statistiques qu'on en déduit »}).
  \item Implémenter fidèlement, sans simplification algébrique
  intermédiaire, les formules dérivées par M4 pour le coût pénalisé, son
  gradient et le solveur, de façon à préserver la correspondance entre le
  rapport et le code (exigence citée dans \texttt{gradient.py}\footnote{Le
  fichier source la nomme « Contrainte méthodologique 1 ».}).
  \item Confronter empiriquement deux résultats théoriques du projet au
  comportement observé du code : la borne de convergence $\eta < 2/L$
  (Expérience 4) et le lien entre maillage du réseau et conditionnement
  (Expérience 5), ainsi qu'analyser la sensibilité du problème au paramètre
  de pénalisation $\mu$ (Expérience 3).
\end{enumerate}

% =====================================================================
\section{Interfaces avec les autres membres}
% =====================================================================

\subsection{Dépendances vis-à-vis de M3}

M5 dépend du module \texttt{src/probability/demand\_model.py}, rédigé par M3,
via les fonctions suivantes, effectivement importées dans
\texttt{monte\_carlo.py} :

\begin{itemize}
  \item \texttt{parametres\_demande(reseau)} $\to (\mu, \sigma)$ : vecteurs
  des moyennes et écarts-types de demande par quartier.
  \item \texttt{matrice\_covariance(reseau)} $\to \Sigma$ : matrice de
  covariance des demandes, construite comme $\Sigma_{ij} = \rho_{ij}
  \sigma_i \sigma_j$ à partir des corrélations déclarées dans
  \texttt{network\_config.json}. Cette fonction lève une exception si
  $\Sigma$ n'est pas définie positive.
  \item \texttt{echantillonner(mu, covariance, n\_tirages, generateur)}
  $\to$ tableau $(n\_tirages, p)$ : tirages gaussiens multivariés, tronqués à
  zéro par en bas (\texttt{np.clip(..., a\_min=0.0, ...)}) pour préserver la
  positivité physique des débits.
  \item \texttt{estimateur\_moyenne}, \texttt{estimateur\_variance},
  \texttt{correlations\_empiriques} : importées par \texttt{monte\_carlo.py}
  mais non appelées dans le corps de ce fichier tel que fourni. Leur usage
  effectif par M5 n'apparaît pas dans les fichiers transmis.
\end{itemize}

M5 ne dépend pas de \texttt{intervalle\_confiance\_moyenne} ni de
\texttt{quartiers\_atypiques}, bien que ces deux fonctions existent dans
\texttt{demand\_model.py} et semblent conçues pour le type d'analyse que
l'Expérience 6 revendique (voir section~15).

Le contrat d'interface entre \texttt{demand\_model.py} et
\texttt{monte\_carlo.py} est cohérent : les six fonctions importées existent
bien, avec les signatures attendues.

\subsection{Dépendances vis-à-vis de M4}

M5 dépend d'un document théorique de M4,
\texttt{docs/section\_optimisation\_membre4.pdf} (non fourni), pour trois
résultats repris littéralement dans le code :

\begin{itemize}
  \item la formulation pénalisée $J(q) = q^\top C q + \mu \lVert Aq - b
  \rVert^2$ (section~7 du docstring de \texttt{objective.py}) ;
  \item le gradient $\nabla J(q) = 2Cq + 2\mu A^\top (Aq - b)$ (section~8) ;
  \item la borne de convergence $\eta < 2/L$ avec $L = \lambda_{\max}(2C +
  2\mu A^\top A)$, ainsi que le choix pratique usuel $\eta = 1/L$.
\end{itemize}

Un écart assumé avec M4 est explicitement documenté dans le code et doit
être arbitré par le groupe : le document de M4 (section~5.3, selon le
docstring de \texttt{gradient\_descent.py}) propose trois critères d'arrêt et
retient par défaut la norme du gradient $\lVert \nabla J \rVert \to 0$. M5
argumente que ce critère n'est pas adapté à un problème projeté : une
composante dont le débit optimal vaut zéro conserve un gradient non nul, que
la projection annule à chaque itération, ce qui empêcherait la boucle de
s'arrêter avant \texttt{max\_iterations}. M5 retient donc par défaut le
critère de déplacement $\lVert q_{k+1} - q_k \rVert < \text{tolérance}$, tout
en rendant les trois critères de M4 disponibles via le paramètre
\texttt{critere} de \texttt{descente\_projetee}. Ce point est signalé dans le
code comme devant être arbitré par M4 pour harmoniser le rapport final.

\subsection{Interfaces avec M6}

Voir la section~15, dédiée spécifiquement à cette interface.

% =====================================================================
\section{Architecture logicielle de ma partie}
% =====================================================================

Quatre modules constituent l'implémentation de M5, avec les dépendances
suivantes entre eux (flèche = « est utilisé par ») :

\begin{verbatim}
demand_model.py (M3)
      |
      v
monte_carlo.py --------------------> [non consommé par les scripts fournis]

objective.py --------> gradient_descent.py --------> run_scenarios.py
      ^                       ^                              |
      |                       |                              v
   exp3, exp4, exp5 appellent directement            [non consommé par
   descente_projetee (donc gradient_descent.py         exp3, exp4, exp5]
   et objective.py), sans passer par run_scenarios.py
\end{verbatim}

Ce schéma reflète un constat d'audit à porter à la connaissance du groupe :
\textbf{aucun des scripts d'expérience fournis (3, 4, 5) n'appelle
\texttt{run\_scenarios.py} ni \texttt{monte\_carlo.py}.} Les trois
expériences disponibles résolvent un problème déterministe unique, à la
demande moyenne $\mu$ obtenue via \texttt{parametres\_demande}, et non sur un
ensemble de scénarios Monte-Carlo. Le module \texttt{monte\_carlo.py} et les
fonctions batch de \texttt{run\_scenarios.py} (\texttt{resoudre\_tous},
\texttt{descente\_projetee\_par\_lot}, \texttt{scenarios\_de\_stress}) sont
implémentés mais non démontrés empiriquement par les fichiers transmis. Ce
point est développé en section~14.

% =====================================================================
\section{Description détaillée des modules}
% =====================================================================

\subsection{\texttt{monte\_carlo.py}}

\begin{description}[style=nextline]
  \item[Rôle] Transformer le modèle probabiliste de demande de M3 en
  scénarios exploitables et en statistiques agrégées.
  \item[Objectif] Ne pas se limiter au tirage de nombres aléatoires : fournir
  une mesure explicite de la précision des statistiques déduites (erreur
  standard, quantiles).
  \item[Entrées] Un objet \texttt{reseau} (ou dictionnaire de configuration),
  un nombre de tirages, une graine optionnelle.
  \item[Sorties] Un tableau de scénarios $(n\_tirages, |\text{quartiers}|)$,
  ou un objet \texttt{StatistiquesMonteCarlo} (moyenne, variance, erreur
  standard, quantiles, nombre de tirages).
  \item[Principales fonctions] \texttt{generer\_scenarios},
  \texttt{agreger}, \texttt{erreur\_estimation},
  \texttt{convergence\_par\_taille}.
  \item[Algorithmes] Tirage gaussien multivarié (délégué à
  \texttt{demand\_model.echantillonner}) ; calcul de moyenne, variance
  corrigée (\texttt{ddof=1}) et quantiles empiriques via \texttt{numpy}.
  \item[Dépendances] \texttt{src.probability.demand\_model} (M3).
  \item[Hypothèses] La graine par défaut (\texttt{GRAINE\_PAR\_DEFAUT = 42})
  est fixée au niveau du module plutôt qu'à chaque appel, pour qu'un oubli de
  graine produise un résultat reproductible plutôt qu'un résultat différent à
  chaque exécution.
  \item[Contrôles] \texttt{agreger} lève une exception si moins de deux
  valeurs sont fournies (variance corrigée indéfinie). \texttt{
  convergence\_par\_taille} lève une exception si une taille demandée dépasse
  la population disponible, plutôt que de réduire silencieusement la taille.
  \item[Fichiers qui l'utilisent] Aucun des fichiers fournis (\texttt{exp3},
  \texttt{exp4}, \texttt{exp5}, \texttt{exp6}, \texttt{run\_scenarios.py})
  n'importe ce module. Son docstring affirme qu'il alimente
  \texttt{run\_scenarios} et les Expériences 2 et 6 ; cette affirmation n'est
  vérifiée par aucun des fichiers transmis (voir section~15).
\end{description}

\subsection{\texttt{objective.py}}

\begin{description}[style=nextline]
  \item[Rôle] Implémenter la fonction de coût pénalisée $J(q)$, son
  gradient, sa hessienne et un contrôle de correction du gradient.
  \item[Objectif] Traduire fidèlement, sans simplification algébrique, les
  formules dérivées par M4.
  \item[Entrées] $q \in \mathbb{R}^{|E|}$ (débits), $A$ (matrice
  d'incidence, forme $(|V|,|E|)$), $b \in \mathbb{R}^{|V|}$ (second membre de
  conservation), $c \in \mathbb{R}^{|E|}$ (coûts par conduite), $\mu > 0$
  (paramètre de pénalisation).
  \item[Sorties] Un scalaire ($J(q)$, coût technique, norme de violation) ou
  un vecteur/matrice (gradient, hessienne).
  \item[Principales fonctions] \texttt{cout}, \texttt{violation\_contrainte},
  \texttt{objectif}, \texttt{gradient}, \texttt{hessienne},
  \texttt{verifier\_gradient}.
  \item[Algorithmes] Calcul direct des formes quadratiques ; différences
  finies centrées d'ordre $O(h^2)$ pour la vérification numérique du
  gradient.
  \item[Dépendances] Aucune dépendance à un autre module du projet ; utilise
  uniquement \texttt{numpy}.
  \item[Hypothèses] Les coûts $c_e$ sont strictement positifs (condition de
  définie-positivité de $C = \mathrm{diag}(c_e)$, nécessaire à la stricte
  convexité de $J$).
  \item[Contrôles] \texttt{verifier\_gradient} compare le gradient
  analytique aux différences finies centrées ; un écart relatif maximal de
  l'ordre de $10^{-6}$ ou moins indique une dérivation correcte. Ce contrôle
  ne remplace pas la dérivation manuscrite exigée par le sujet, il la
  vérifie.
  \item[Fichiers qui l'utilisent] \texttt{gradient\_descent.py} (import
  direct de \texttt{gradient}, \texttt{hessienne}, \texttt{objectif},
  \texttt{violation\_contrainte}) ; \texttt{run\_scenarios.py} (\texttt{cout},
  \texttt{violation\_contrainte}) ; \texttt{exp3\_sensitivity\_mu.py}
  (\texttt{cout}, \texttt{hessienne}, \texttt{violation\_contrainte}) ;
  \texttt{exp5\_conditioning\_vs\_mesh.py} (\texttt{hessienne}).
\end{description}

\textbf{Formulation mathématique implémentée} (dérivée par M4, section 6 à 8
de son document) :

\begin{align}
  \text{Problème d'origine :} \quad & q^* = \arg\min_q \sum_{e \in E} c_e
  q_e^2 \quad \text{sous } Aq = b,\ q \geq 0 \label{eq:probleme} \\
  \text{Pénalisation :} \quad & J(q) = q^\top C q + \mu \lVert Aq - b
  \rVert^2 \label{eq:objectif} \\
  \text{Gradient :} \quad & \nabla J(q) = 2Cq + 2\mu A^\top(Aq - b)
  \label{eq:gradient}
\end{align}

Le module ne redérive pas la démonstration de convexité de M4 : il en
implémente la conséquence directe (voir 6.3 pour la hessienne constante $H =
2C + 2\mu A^\top A$).

\subsection{\texttt{gradient\_descent.py}}

\begin{description}[style=nextline]
  \item[Rôle] Résoudre le problème pénalisé~\eqref{eq:objectif} par descente
  de gradient projetée sur l'orthant positif.
  \item[Objectif] Fournir un solveur robuste, avec traçabilité de la
  convergence et garde-fou contre la divergence numérique.
  \item[Entrées] $A$, $b$, $c$, $\mu$, pas $\eta$ (optionnel), point de
  départ $q_0$ (optionnel), tolérance, nombre maximal d'itérations, critère
  d'arrêt (\texttt{"deplacement"}, \texttt{"gradient"} ou \texttt{"cout"}).
  \item[Sorties] Le couple $(q^*, \text{historique})$, où
  \texttt{HistoriqueConvergence} trace les valeurs de $J(q_k)$, les normes de
  gradient, les violations de contrainte, le nombre d'itérations, le
  statut de convergence et le motif d'arrêt.
  \item[Principales fonctions] \texttt{projeter\_orthant\_positif},
  \texttt{pas\_maximal\_theorique}, \texttt{descente\_projetee}.
  \item[Algorithmes] Voir pseudo-code ci-dessous. Projection euclidienne sur
  $\{q \geq 0\}$ par $P(x) = \max(x, 0)$, appliquée à \emph{chaque}
  itération (et non une seule fois à la fin, erreur explicitement signalée
  comme un autre algorithme qui ne converge pas vers le même point).
  \item[Dépendances] \texttt{objective.py}.
  \item[Hypothèses] Si $\eta$ n'est pas fourni, il est fixé à la moitié de la
  borne théorique $2/L$, soit $\eta = 1/L$, et ce choix doit être signalé
  dans le rapport plutôt que présenté comme un réglage manuel. Si $q_0$
  n'est pas fourni, le point de départ est le vecteur nul, admissible pour
  $q \geq 0$.
  \item[Contrôles] Un seuil de divergence (\texttt{SEUIL\_DIVERGENCE = 1e12})
  arrête proprement la boucle si le déplacement devient non fini ou
  excessif, plutôt que de laisser la sortie se remplir de \texttt{inf}/
  \texttt{nan}. Ce seuil est un garde-fou de mise en \oe uvre introduit par
  M5, non issu du document de M4.
  \item[Fichiers qui l'utilisent] \texttt{run\_scenarios.py}
  (\texttt{descente\_projetee}, \texttt{pas\_maximal\_theorique}) ;
  \texttt{exp3\_sensitivity\_mu.py}, \texttt{exp4\_convergence\_check.py},
  \texttt{exp5\_conditioning\_vs\_mesh.py} (\texttt{descente\_projetee},
  \texttt{pas\_maximal\_theorique}).
\end{description}

\textbf{Pseudo-code} (tel que documenté dans le module) :

\begin{quote}
\textbf{Entrées} : $A, b, c, \mu, \eta, q_0$, tolérance, $k_{\max}$\\
\textbf{Pour} $k = 0, 1, 2, \dots$ :
\begin{itemize}
  \item $g_k \leftarrow \nabla J(q_k) = 2Cq_k + 2\mu A^\top(Aq_k - b)$
  \item $q_{k+1} \leftarrow P(q_k - \eta g_k)$, avec $P(x) = \max(x, 0)$
  \item arrêt si $\lVert q_{k+1} - q_k \rVert < \text{tolérance}$ ou $k =
  k_{\max}$
\end{itemize}
\textbf{Sortie} : $q_k$
\end{quote}

Deux justifications explicitement portées par le code :

\begin{enumerate}
  \item $\max(\cdot, 0)$ est la projection euclidienne sur $\{q \geq 0\}$
  car le problème $\min_{y \geq 0} \lVert y - x \rVert^2$ se sépare
  composante par composante ; cette séparabilité ne tiendrait pas sur un
  ensemble couplant les composantes (par exemple des contraintes de
  capacité $q_e \leq \text{cap}_e$).
  \item $\eta < 2/L$ garantit la convergence, avec $L = \lambda_{\max}(2C +
  2\mu A^\top A)$ : résultat classique de la descente de gradient projetée
  sur fonction convexe à gradient $L$-lipschitzien. Cette borne théorique
  est confrontée à l'expérience dans l'Expérience 4 (section~10.2).
\end{enumerate}

\subsection{\texttt{run\_scenarios.py}}

\begin{description}[style=nextline]
  \item[Rôle] Résoudre le problème d'optimisation sur un ensemble de
  scénarios de demande, et non uniquement sur la demande moyenne.
  \item[Objectif] Industrialiser le test de robustesse de $q^*$ face à des
  situations variées, exigence explicite du sujet.
  \item[Entrées] Un \texttt{reseau}, une matrice d'incidence $A$, un ou
  plusieurs vecteurs de demande, $\mu$, options du solveur.
  \item[Sorties] Un \texttt{ResultatScenario} (indice, demandes, $q^*$, coût
  technique, violation, nombre d'itérations, convergence) ou une liste de
  ces objets.
  \item[Principales fonctions] \texttt{resoudre\_scenario},
  \texttt{resoudre\_tous}, \texttt{descente\_projetee\_par\_lot},
  \texttt{scenarios\_de\_stress}.
  \item[Algorithmes] \texttt{descente\_projetee\_par\_lot} réplique
  exactement l'algorithme de \texttt{descente\_projetee} (même pas, même
  règle de mise à jour, même projection à chaque itération, même critère de
  déplacement), mais fait avancer $N$ scénarios simultanément dans des
  tableaux de forme $(N, |E|)$ au lieu d'une boucle Python scénario par
  scénario. Le docstring justifie ce choix par un gain de performance
  (environ 8000 itérations nécessaires à $\mu = 100$ quel que soit le nombre
  de scénarios ; en boucle, cela représenterait de l'ordre d'un quart
  d'heure pour 1000 scénarios).
  \item[Dépendances] \texttt{src.graph.build\_graph}
  (\texttt{construire\_second\_membre}, \texttt{vecteur\_couts} --- M1/M2,
  non fournis), \texttt{src.optimization.gradient\_descent},
  \texttt{src.optimization.objective}, \texttt{src.probability.demand\_model}
  (\texttt{parametres\_demande}).
  \item[Hypothèses] Un scénario ayant atteint son seuil de convergence
  continue d'être mis à jour jusqu'à la fin de la boucle par lot ; le
  docstring affirme que cela ne le déplace pas, car il est au point fixe de
  l'itération projetée.
  \item[Contrôles] Le docstring mentionne un test
  \texttt{test\_le\_lot\_donne\_le\_meme\_resultat\_que\_la\_boucle}
  comparant les deux implémentations (boucle et lot) sur le réseau réel.
  \textbf{Ce fichier de test n'a pas été transmis} ; son existence et son
  contenu réel restent à vérifier. Information non disponible dans les
  fichiers fournis.
  \item[Fichiers qui l'utilisent] Aucun des scripts d'expérience fournis
  (\texttt{exp3}, \texttt{exp4}, \texttt{exp5}, \texttt{exp6}) n'importe ce
  module.
\end{description}

\subsection{Scripts expérimentaux}

Voir section~10 pour le détail complet. Le tableau~\ref{tab:experiences}
résume le statut de chaque script fourni.

\begin{longtable}{L{2.2cm}L{2.5cm}L{4.5cm}L{4cm}}
\caption{Statut des scripts expérimentaux fournis.}
\label{tab:experiences}\\
\toprule
Script & Responsable déclaré & Modules M5 utilisés & Remarque \\
\midrule
\endhead
\texttt{exp3\_sensitivity\_mu.py} & M5 (interprétation M4) & \texttt{objective.py}, \texttt{gradient\_descent.py} & Demande moyenne uniquement, pas de Monte-Carlo \\
\texttt{exp4\_convergence\_check.py} & M5 (interprétation M4) & \texttt{objective.py} (via \texttt{gradient\_descent.py}), \texttt{gradient\_descent.py} & idem \\
\texttt{exp5\_conditioning\_vs\_mesh.py} & M5 (interprétation M2) & \texttt{objective.py}, \texttt{gradient\_descent.py} & idem \\
\texttt{exp6\_at\_risk\_districts.py} & M3, selon l'en-tête du fichier lui-même & Aucun & Voir section~15 \\
\bottomrule
\end{longtable}

% =====================================================================
\section{Implémentation de la simulation Monte-Carlo}
% =====================================================================

Cette section documente le module \texttt{monte\_carlo.py} tel qu'implémenté.
Elle doit être lue avec la réserve posée en section~6.1 : aucun des scripts
fournis n'exerce ce module.

\begin{description}[style=nextline]
  \item[Génération des scénarios] \texttt{generer\_scenarios(reseau,
  n\_tirages, graine=None)} récupère $(\mu, \_)$ via
  \texttt{parametres\_demande}, la matrice de covariance $\Sigma$ via
  \texttt{matrice\_covariance}, instancie un générateur \texttt{numpy}
  (graine 42 par défaut) et délègue le tirage à
  \texttt{demand\_model.echantillonner}.
  \item[Utilisation du modèle probabiliste] Le modèle utilisé est
  entièrement celui de M3 : $D \sim \mathcal{N}(\mu, \Sigma)$, tronqué à zéro
  par en bas après tirage.
  \item[Nombre de scénarios] Non fixé dans le code de M5 : c'est un
  paramètre (\texttt{n\_tirages}) laissé à l'appelant. Aucun des fichiers
  fournis n'appelle \texttt{generer\_scenarios} avec une valeur concrète.
  Information non disponible dans les fichiers fournis quant au nombre de
  scénarios effectivement retenu pour la suite du projet.
  \item[Structure d'un scénario] Une ligne d'un tableau $(n\_tirages,
  |\text{quartiers}|)$, chaque colonne correspondant à un quartier, dans
  l'ordre retourné par \texttt{demand\_model}.
  \item[Transmission au solveur] Non exercée dans les fichiers fournis :
  \texttt{run\_scenarios.resoudre\_tous} attend un tableau
  \texttt{scenarios} \emph{issu de} \texttt{generer\_scenarios} selon son
  docstring, mais aucun appel effectif reliant les deux modules n'apparaît
  dans le code transmis.
  \item[Agrégation des résultats] \texttt{agreger(valeurs, quantiles=(0.05,
  0.25, 0.5, 0.75, 0.95))} calcule moyenne, variance corrigée (\texttt{ddof=1}),
  erreur standard $\hat{\sigma}/\sqrt{N}$ et les quantiles demandés.
  \item[Statistiques calculées] Moyenne, variance, erreur standard,
  quantiles aux niveaux 0,05/0,25/0,5/0,75/0,95 par défaut.
  \item[Intervalles de confiance] \textbf{Non implémentés dans
  \texttt{monte\_carlo.py}.} Le module calcule une erreur standard, pas un
  intervalle de confiance au sens statistique (borne de Student ou
  gaussienne). La fonction \texttt{intervalle\_confiance\_moyenne}, qui
  calcule un IC de Student à $(n-1)$ degrés de liberté, existe dans
  \texttt{demand\_model.py} (M3) mais n'est pas importée par
  \texttt{monte\_carlo.py}.
\end{description}

La formule d'erreur standard implémentée,
\begin{equation}
  \text{erreur\_estimation}(\hat{\sigma}, N) = \frac{\hat{\sigma}}{\sqrt{N}},
\end{equation}
est qualifiée dans le code de « volontairement triviale et isolée », le
docstring précisant qu'elle est destinée à être confrontée à la décroissance
mesurée quand $N$ augmente par l'Expérience 2. \textbf{Cette confrontation
n'apparaît dans aucun des fichiers fournis pour l'Expérience 2}, qui dépend
du travail de M6 selon les informations transmises initialement. Une
confrontation empirique de la décroissance en $1/\sqrt{N}$ existe bien dans
les fichiers fournis, mais dans \texttt{exp6\_at\_risk\_districts.py}
(Figure~3, attribuée à M3), pas dans une Expérience 2 attribuée à M5 --- voir
section~15.

\texttt{convergence\_par\_taille(valeurs, tailles=(100, 1000, 10000))}
recalcule les statistiques sur des \emph{préfixes croissants} de la même
population de tirages (et non sur des tirages indépendants à chaque taille),
choix motivé dans le code par la volonté d'observer une même expérience
s'affiner plutôt que de comparer trois expériences différentes.

% =====================================================================
\section{Implémentation du solveur}
% =====================================================================

\begin{description}[style=nextline]
  \item[Fonction objectif] $J(q) = q^\top C q + \mu \lVert Aq - b
  \rVert^2$, avec $C = \mathrm{diag}(c_e)$ (équation~\eqref{eq:objectif}).
  \item[Gradient] $\nabla J(q) = 2Cq + 2\mu A^\top(Aq - b)$
  (équation~\eqref{eq:gradient}), implémenté littéralement d'après la
  dérivation de M4, sans réécriture algébrique intermédiaire. Une note de
  performance sans effet sur le résultat est appliquée : le produit est
  calculé comme $A^\top(Aq - b)$ (deux produits matrice-vecteur) plutôt que
  $(A^\top A)q - A^\top b$ (qui construirait une matrice $|E| \times |E|$ à
  chaque appel).
  \item[Pénalisation] Le paramètre $\mu$ contrôle le compromis entre respect
  de la conservation des flux et coût technique. La solution pénalisée ne
  satisfait $Aq = b$ qu'à la limite $\mu \to \infty$ ; le résidu $\lVert Aq^*
  - b \rVert$ est une quantité à mesurer et non à supposer nulle.
  \item[Descente de gradient] Voir section~6.3 pour le pseudo-code complet.
  \item[Projection] Projection euclidienne sur $\{q \geq 0\}$ à chaque
  itération, $P(x) = \max(x, 0)$.
  \item[Critère d'arrêt] Par défaut, critère de déplacement $\lVert q_{k+1}
  - q_k \rVert < \text{tolérance}$ (choix de M5, motivé et documenté en
  section~4.2). Les critères de norme de gradient et de variation de coût
  sont également disponibles via le paramètre \texttt{critere}.
  \item[Gestion des paramètres] $\eta$ par défaut : $\tfrac{1}{2} \times
  \tfrac{2}{L}= \tfrac{1}{L}$. $q_0$ par défaut : vecteur nul. Tolérance par
  défaut : $10^{-8}$. \texttt{max\_iterations} par défaut : $10\,000$ pour
  \texttt{descente\_projetee}, $100\,000$ pour
  \texttt{descente\_projetee\_par\_lot}.
  \item[Suivi de la convergence] \texttt{HistoriqueConvergence} enregistre,
  à chaque itération, la valeur de $J(q_k)$, la norme du gradient et la
  violation de contrainte, pour éviter d'avoir à relancer l'optimisation
  pour tracer une courbe a posteriori.
  \item[Résultats retournés] Le couple $(q^*, \text{historique})$, avec dans
  l'historique le nombre d'itérations effectuées, un booléen de convergence
  et le motif d'arrêt (texte libre : critère atteint, divergence détectée,
  ou nombre maximal d'itérations atteint).
\end{description}

Un garde-fou de divergence (\texttt{SEUIL\_DIVERGENCE = 1e12}) interrompt la
boucle proprement si le déplacement devient non fini ou dépasse ce seuil ---
utile en particulier pour l'Expérience 4, qui dépasse volontairement la borne
théorique $2/L$.

% =====================================================================
\section{Orchestration des scénarios}
% =====================================================================

Le pipeline décrit dans \texttt{run\_scenarios.py} est le suivant :

\begin{center}
demande simulée $\to$ construction du second membre $b$ (via
\texttt{construire\_second\_membre}, M1/M2) $\to$ optimisation
(\texttt{descente\_projetee} ou \texttt{descente\_projetee\_par\_lot}) $\to$
solution $q^*$ $\to$ collecte des résultats (\texttt{ResultatScenario}) $\to$
[agrégation --- non implémentée dans ce module] $\to$ analyse expérimentale.
\end{center}

Deux réserves à porter à la connaissance du groupe :

\begin{enumerate}
  \item \textbf{L'étape d'agrégation n'est pas câblée.} Ni
  \texttt{resoudre\_scenario} ni \texttt{resoudre\_tous} n'appellent
  \texttt{monte\_carlo.agreger} sur la liste de résultats produite. La
  capacité d'agrégation statistique existe (section~7) mais n'est pas
  reliée à l'orchestration des scénarios dans les fichiers fournis.
  \item \textbf{Ce module n'est exercé par aucune expérience fournie.} Les
  fonctions \texttt{resoudre\_scenario}, \texttt{resoudre\_tous},
  \texttt{descente\_projetee\_par\_lot} et \texttt{scenarios\_de\_stress}
  sont implémentées mais non appelées par \texttt{exp3}, \texttt{exp4},
  \texttt{exp5} ou \texttt{exp6}. Elles sont vraisemblablement destinées aux
  Expériences 1 et 2, dépendantes du travail de M6 selon les informations
  transmises.
\end{enumerate}

\texttt{scenarios\_de\_stress(reseau, quantile=0.95)} construit un scénario
de demande \emph{extrême marginal} : chaque quartier est individuellement
poussé à son propre quantile ($\mu + z_{0.95} \sigma$, composante par
composante), ce qui diffère du quantile de la demande totale du réseau
(généralement plus faible, car les écarts individuels ne se cumulent jamais
parfaitement). Le code signale explicitement que le rapport final doit
préciser laquelle des deux lectures est retenue --- ce choix n'est pas encore
tranché dans les fichiers fournis.

% =====================================================================
\section{Description des expériences}
% =====================================================================

\subsection{Expérience 3 --- Sensibilité au paramètre de pénalisation
$\mu$}

\begin{description}[style=nextline]
  \item[Objectif] Mesurer le compromis porté par $\mu$ entre respect de la
  contrainte de conservation et vitesse de convergence.
  \item[Question scientifique] Comment le choix de $\mu$ affecte-t-il
  simultanément la violation résiduelle et le nombre d'itérations
  nécessaires à la convergence ?
  \item[Paramètres contrôlés] $\mu \in \{1, 10, 100, 1000\}$ ; tolérance
  $10^{-8}$ ; \texttt{max\_iterations} $= 200\,000$.
  \item[Méthode] Pour chaque $\mu$, résoudre le problème pénalisé à la
  demande moyenne avec \texttt{descente\_projetee}, en recalculant $\eta =
  \tfrac{1}{2} \times \tfrac{2}{L}$ à chaque valeur de $\mu$ (protocole
  explicitement retenu : le nombre d'itérations mesuré mélange donc l'effet
  du durcissement du problème et celui du rétrécissement du pas, ce qui est
  volontaire car les deux effets partagent la même origine algébrique, la
  matrice $2C + 2\mu A^\top A$).
  \item[Protocole expérimental] Réseau synthétique, demande moyenne
  uniquement (pas de tirage Monte-Carlo).
  \item[Métriques] Valeur propre maximale $L$, conditionnement de $H$, pas
  maximal théorique $2/L$, pas utilisé, nombre d'itérations, convergence,
  coût technique, violation résiduelle.
  \item[Sorties produites] Table CSV \texttt{results/tables/
  exp3\_sensibilite\_mu.csv} (chemin obtenu via \texttt{\_bootstrap.TABLEAUX},
  module non fourni).
  \item[Figures] \texttt{results/figures/exp3\_sensibilite\_mu.png}, deux
  panneaux : à gauche, $J(q_k)$ en échelle log-log en fonction de $k$, une
  courbe par $\mu$ ; à droite, violation résiduelle et nombre d'itérations en
  fonction de $\mu$, en échelle log-log.
  \item[Interprétation attendue] La violation résiduelle décroît quand $\mu$
  augmente, mais la convergence ralentit --- lien direct avec la dérivation
  de l'Étape~7 : augmenter $\mu$ resserre la contrainte d'égalité mais
  alourdit $\lambda_{\max}(H)$, dégrade donc $\kappa(H)$, impose donc un pas
  plus petit via $\eta < 2/L$, ralentit donc la convergence.
  \item[Lien avec la théorie] Satisfait l'exigence d'analyse de sensibilité
  d'au moins un hyperparamètre (« Contrainte méthodologique 3, exigence
  (v) », selon le docstring du script).
  \item[Limites] Les valeurs numériques réelles (colonnes du CSV, valeurs de
  $L$, de $\kappa(H)$, nombre d'itérations effectif) ne figurent pas dans le
  code source et n'ont pas été transmises. Information non disponible dans
  les fichiers fournis.
\end{description}

\subsection{Expérience 4 --- Vérification de la borne de convergence}

\begin{description}[style=nextline]
  \item[Objectif] Confronter la borne théorique $\eta < 2/L$ au comportement
  observé.
  \item[Question scientifique] La convergence est-elle stable sous la borne
  et instable au-delà, comme le prédit la théorie ?
  \item[Paramètres contrôlés] $\mu = 100$ ; facteurs de la borne $\eta/(2/L)
  \in \{0{,}1;\ 0{,}5;\ 0{,}9;\ 1{,}1;\ 1{,}5\}$ ; \texttt{max\_iterations} $=
  60\,000$.
  \item[Méthode] Pour chaque facteur, calculer $\eta = \text{facteur}
  \times (2/L)$ et lancer \texttt{descente\_projetee}, en suivant $J(q_k)$ et
  la norme du gradient au fil des itérations.
  \item[Protocole expérimental] Réseau synthétique, demande moyenne, cinq
  valeurs de $\eta$ encadrant franchement le seuil théorique (trois en
  dessous, deux au-dessus).
  \item[Métriques] $J(q_k)$ et $\lVert \nabla J(q_k) \rVert$ en fonction de
  $k$ ; remontée maximale de $J$ entre deux itérations consécutives ;
  convergence atteinte ou non ; motif d'arrêt.
  \item[Sorties produites] Table CSV \texttt{results/tables/
  exp4\_convergence.csv}.
  \item[Figures] \texttt{results/figures/exp4\_convergence.png} : $J(q_k)$
  en échelle logarithmique, une courbe par facteur (trait plein sous le
  seuil, tireté au-dessus), limitée aux 2000 premières itérations affichées.
  \item[Interprétation attendue] Convergence stable pour les facteurs
  inférieurs à 1, oscillation ou divergence au-delà --- confirmée ou non par
  les données numériques réelles, qui n'ont pas été transmises.
  \item[Lien avec la théorie] Confronte le résultat classique de convergence
  de la descente de gradient projetée sur fonction convexe à gradient
  $L$-lipschitzien, et la constante de Lipschitz calculée à partir du $L$
  fourni par M2 (Étape~3).
  \item[Limites] L'erreur à ne pas commettre, explicitement signalée dans le
  docstring, est de présenter une courbe qui descend comme preuve de
  convergence : ce que l'expérience démontre est le lien entre la constante
  de Lipschitz et le comportement de part et d'autre du seuil, pas le simple
  fait que $J$ décroisse. Les valeurs numériques réelles (borne calculée,
  nombre d'itérations par facteur) n'ont pas été transmises. Information non
  disponible dans les fichiers fournis.
\end{description}

\subsection{Expérience 5 --- Maillage du réseau et conditionnement}

\begin{description}[style=nextline]
  \item[Objectif] Vérifier empiriquement le lien entre structure du graphe
  et stabilité numérique établi à l'Étape~3 (M2).
  \item[Question scientifique] Un réseau mal maillé se traduit-il par un
  conditionnement plus élevé et davantage d'itérations nécessaires à la
  convergence ?
  \item[Paramètres contrôlés] Densités de maillage $\{0;\ 0{,}25;\ 0{,}5;\
  0{,}75;\ 1\}$ ; $\mu = 100$ ; tolérance $10^{-8}$ ; \texttt{max\_iterations}
  $= 200\,000$.
  \item[Méthode] Pour chaque variante de maillage (obtenue via
  \texttt{generer\_variantes\_de\_maillage}, M1/M2, non fourni), calculer
  $\kappa(A)$ et $\kappa(H) = \kappa(2C + 2\mu A^\top A)$, résoudre avec un
  pas commun à toutes les variantes (protocole retenu explicitement : $\eta$
  fixé à la moitié de la plus petite borne théorique rencontrée parmi toutes
  les variantes, pour ne pas mélanger l'effet du conditionnement avec celui
  d'un pas recalculé à chaque variante).
  \item[Protocole expérimental] Réseaux synthétiques à densités de connexion
  croissantes, tous vérifiés connexes.
  \item[Métriques] Nombre de conduites, connexité, rang de $A$, dimension du
  noyau (nombre de boucles), $\kappa(A)$, $\kappa(H)$, nombre d'itérations,
  convergence.
  \item[Sorties produites] Table CSV \texttt{results/tables/
  exp5\_maillage\_conditionnement.csv}.
  \item[Figures] \texttt{results/figures/exp5\_maillage\_conditionnement.png},
  deux panneaux : à gauche, $\kappa(A)$ et $\kappa(H)$ en fonction du nombre
  de conduites (axes doubles) ; à droite, nombre d'itérations en fonction de
  $\kappa(H)$, annoté du nombre de boucles par variante.
  \item[Interprétation attendue] Relation croissante (non nécessairement
  linéaire) entre conditionnement et nombre d'itérations. Le docstring
  signale explicitement de ne pas sur-interpréter une régression linéaire
  sur cinq points.
  \item[Lien avec la théorie] C'est $\kappa(H)$, et non $\kappa(A)$, qui
  gouverne la vitesse de la descente de gradient, selon la section~6.3 du
  document de M4 citée dans le code.
  \item[Limites] Conditions de validité rappelées dans le docstring : toutes
  les variantes doivent rester connexes (sinon le rang de $A$ change de
  nature et les $\kappa$ ne sont plus comparables) ; $\mu$, la tolérance
  doivent rester identiques d'une variante à l'autre. Les valeurs numériques
  réelles n'ont pas été transmises. Information non disponible dans les
  fichiers fournis.
\end{description}

\subsection{Expérience 6 --- Quartiers à risque}

Voir la réserve d'attribution détaillée en section~15. Le contenu est
documenté ici tel que transmis.

\begin{description}[style=nextline]
  \item[Objectif déclaré (docstring)] Repérer les quartiers dont la demande
  s'écarte régulièrement des prévisions, et relier ce constat à la
  robustesse de $q^*$.
  \item[Attribution réelle du code] Le fichier s'auto-attribue explicitement
  à M3 dans un bloc de commentaires précédant le code exécutable : « Ce
  script (Étape~4) [\ldots] Implémentation du Membre 3 (Probabilités et
  Statistiques) ». Le code n'importe que
  \texttt{src.probability.demand\_model.DemandModel} ; aucun des quatre
  modules de M5 n'est utilisé.
  \item[Méthode réellement implémentée] Tirage de $N = 10\,000$ scénarios
  via \texttt{DemandModel.sample\_demands} (graine 42), calcul des
  quantiles empiriques à 95~\% et 99~\%, et du pic maximal observé, par
  quartier. Des libellés de type de risque (« Topologique (Pont unique) »
  pour Q1, « Stochastique (Forte variance) » pour Q6, « Transit inter-zones »
  pour Q4/Q5, « Modéré / Standard » sinon) sont codés en dur dans le script,
  sans calcul topologique réel à l'appui : le réseau (objet
  \texttt{Reseau}/\texttt{network\_config.json} via \texttt{charger\_reseau})
  n'est jamais chargé dans ce script, contrairement aux Expériences 3, 4 et
  5.
  \item[Écart avec la méthode annoncée] La fonction
  \texttt{quartiers\_atypiques} de \texttt{demand\_model.py}, qui calcule
  précisément une fréquence de dépassement d'intervalle de confiance par
  quartier --- exactement ce que décrit le docstring de l'Expérience 6 ---
  n'est pas appelée par le script fourni. Le script utilise à la place une
  analyse par quantiles simples (95~\%, 99~\%, maximum observé), sans
  notion de biais volontaire introduit sur des quartiers connus, ni de
  comparaison au taux de fausses alertes attendu de 5~\% mentionné dans le
  docstring.
  \item[Sorties produites] Quatre figures dans \texttt{results/figures/} :
  \begin{enumerate}
    \item \texttt{fig1\_matrices\_correlation\_covariance.png} : matrices de
    corrélation $R$ et de covariance $\Sigma$ du modèle.
    \item \texttt{fig2\_densites\_probabilites\_demande.png} : densités
    gaussiennes de demande par quartier.
    \item \texttt{fig3\_convergence\_loi\_grands\_nombres.png} : erreur
    empirique $\lVert \hat\mu_N - \mu \rVert_2$ en fonction de $N \in
    \{10, 25, 50, 100, 250, 500, 1000, 2500, 5000, 10000, 50000\}$, comparée
    à la pente théorique en $O(1/\sqrt{N})$, en échelle log-log.
    \item \texttt{fig4\_profil\_risque\_quartiers.png} : boîtes à moustaches
    de la demande par quartier, avec un seuil de charge critique locale
    codé en dur à 60~m$^3$/h.
  \end{enumerate}
  \item[Interprétation] Non disponible pour les valeurs numériques ; les
  figures ne sont pas fournies en tant qu'images dans ce document.
  \item[Lien avec la théorie] La Figure~3 réalise concrètement la
  vérification empirique de la décroissance en $1/\sqrt{N}$ que le docstring
  de \texttt{erreur\_estimation} (module de M5) attribue explicitement au
  rôle de l'Expérience 2. Ce recouvrement est à faire trancher par le
  groupe : soit l'Expérience 2 (attribuée à M5, dépendante de M6) est
  redondante avec cette figure, soit elle doit couvrir un aspect différent
  non encore précisé.
  \item[Limites] Le script contient une note d'intégration indiquant qu'un
  squelette d'origine, qui levait une \texttt{NotImplementedError}, faisait
  planter le script avant d'atteindre le code utile ; les quatre figures
  ont donc été produites en appelant \texttt{run\_experiment\_6} directement
  plutôt que via le point d'entrée \texttt{main()} tel qu'exécuté
  normalement. Le fichier fourni corrige cet ordre, mais ce point historique
  est à connaître pour la reproductibilité (section~12).
\end{description}

% =====================================================================
\section{Résultats obtenus}
% =====================================================================

Aucun fichier de résultats (CSV de \texttt{results/tables/}, images de
\texttt{results/figures/}) n'a été transmis avec cette demande. Les scripts
fournis produisent du code capable de générer ces résultats, mais
\textbf{aucune valeur numérique réelle} (conditionnement mesuré, nombre
d'itérations effectif, violation résiduelle observée, etc.) n'apparaît dans
les fichiers source eux-mêmes : ces valeurs ne sont connues qu'à l'exécution.

Information non disponible dans les fichiers fournis pour l'ensemble de
cette section. Le tableau~\ref{tab:resultats-attendus} récapitule uniquement
la nature des résultats que chaque expérience \emph{produirait} à
l'exécution, sans préjuger de leur valeur.

\begin{longtable}{L{2cm}L{2.7cm}L{4.5cm}L{4cm}}
\caption{Nature des résultats attendus par expérience (valeurs non
disponibles).}
\label{tab:resultats-attendus}\\
\toprule
Expérience & Fichier source & Table attendue & Figure attendue \\
\midrule
\endhead
3 & \texttt{results/tables/ exp3\_sensibilite\_mu.csv} & $\mu$, $L$,
$\kappa(H)$, pas, itérations, coût, violation & \texttt{exp3\_sensibilite\_mu.png} \\
4 & \texttt{results/tables/ exp4\_convergence.csv} & facteur de la borne,
$\eta$, convergence, itérations, $J$ final & \texttt{exp4\_convergence.png} \\
5 & \texttt{results/tables/ exp5\_maillage\_conditionnement.csv} & densité,
conduites, $\kappa(A)$, $\kappa(H)$, itérations & \texttt{exp5\_maillage\_conditionnement.png} \\
6 & --- (attribué à M3) & non produite (impression console uniquement) & 4 figures listées en section~10.4 \\
\bottomrule
\end{longtable}

Pour compléter cette section, transmettre à M5 (ou directement à M6) le
contenu réel de \texttt{results/tables/} et, si possible, les images de
\texttt{results/figures/}.

% =====================================================================
\section{Reproductibilité}
% =====================================================================

\begin{description}[style=nextline]
  \item[Environnement] Python avec \texttt{numpy}, \texttt{scipy}
  (\texttt{scipy.stats}), \texttt{matplotlib} (backend \texttt{Agg} imposé
  explicitement dans les scripts d'expérience, pour un rendu sans affichage
  interactif). Version exacte de Python et des dépendances : information non
  disponible dans les fichiers fournis (pas de \texttt{requirements.txt} ni
  de fichier d'environnement transmis).
  \item[Commandes] Chaque script d'expérience s'exécute indépendamment,
  via son bloc \texttt{if \_\_name\_\_ == "\_\_main\_\_":} qui appelle
  \texttt{main()}. Exemple : \texttt{python experiments/
  exp3\_sensitivity\_mu.py}. Un module \texttt{\_bootstrap} est importé en
  tête de chaque script d'expérience (\texttt{3}, \texttt{4}, \texttt{5}) et
  expose au minimum \texttt{CONFIG}, \texttt{TABLEAUX} et \texttt{FIGURES} ;
  son contenu réel n'a pas été transmis. Information non disponible dans les
  fichiers fournis.
  \item[Ordre d'exécution] Aucune dépendance d'ordre n'existe entre les
  Expériences 3, 4 et 5 : chacune recharge le réseau et reconstruit ses
  propres objets ($A$, $b$, coûts) indépendamment. L'Expérience 6, telle que
  fournie, ne dépend d'aucune des trois autres.
  \item[Dépendances] \texttt{exp3}, \texttt{exp4} et \texttt{exp5} dépendent
  de \texttt{data.generate\_network.charger\_reseau} (M1),
  \texttt{src.graph.build\_graph} (M1/M2) et, pour \texttt{exp5},
  \texttt{data.generate\_network.generer\_variantes\_de\_maillage} et
  \texttt{src.graph.graph\_analysis} (M2). \texttt{exp6} dépend uniquement
  de \texttt{src.probability.demand\_model.DemandModel}, qui lit
  \texttt{data/network\_config.json} (avec un repli sur un chemin relatif
  \texttt{../../data/network\_config.json} si le premier chemin n'existe
  pas) et retourne un dictionnaire vide si aucun des deux fichiers n'est
  trouvé --- ce qui ferait alors échouer silencieusement le calcul de
  $\mu, \sigma, \Sigma$ dans \texttt{parametres\_demande}/
  \texttt{matrice\_covariance} plutôt que de signaler l'absence de fichier.
  \item[Fichiers générés] \texttt{results/tables/*.csv} et
  \texttt{results/figures/*.png}, selon le détail donné en section~10 et
  résumé au tableau~\ref{tab:resultats-attendus}.
\end{description}

% =====================================================================
\section{Tests et validation de l'implémentation}
% =====================================================================

Aucun fichier de test n'a été transmis avec cette demande. Les fichiers
source de M5 \textbf{référencent} l'existence de tests, sans que leur
contenu soit disponible :

\begin{itemize}
  \item \texttt{tests/test\_projection.py}, mentionné dans le docstring de
  \texttt{projeter\_orthant\_positif} comme testant cette fonction de
  manière isolée.
  \item Un test nommé
  \texttt{test\_le\_lot\_donne\_le\_meme\_resultat\_que\_la\_boucle},
  mentionné dans le docstring de \texttt{descente\_projetee\_par\_lot} comme
  comparant les résultats de la version boucle et de la version lot sur le
  réseau réel.
\end{itemize}

Le seul contrôle numérique dont le code source complet est disponible est
\texttt{verifier\_gradient} (\texttt{objective.py}), qui compare le gradient
analytique aux différences finies centrées et retourne l'écart relatif
maximal, avec un seuil indicatif de $10^{-6}$ pour une dérivation jugée
correcte. Il s'agit d'une fonction de contrôle, pas d'un test automatisé au
sens \texttt{pytest} : rien dans les fichiers fournis n'indique qu'elle soit
appelée par une suite de tests.

Information non disponible dans les fichiers fournis : contenu réel des
tests unitaires, tests d'intégration, couverture de code, cas limites
couverts. \textbf{À demander explicitement au groupe} pour compléter cette
section.

% =====================================================================
\section{Limites et points de vigilance}
% =====================================================================

\begin{enumerate}
  \item \textbf{Le module \texttt{monte\_carlo.py} n'est exercé par aucune
  expérience fournie.} Ses fonctions de génération de scénarios et
  d'agrégation statistique sont implémentées mais non démontrées
  empiriquement dans les fichiers transmis.
  \item \textbf{Les fonctions batch de \texttt{run\_scenarios.py} ne sont
  exercées par aucune expérience fournie.} \texttt{resoudre\_tous},
  \texttt{descente\_projetee\_par\_lot} et \texttt{scenarios\_de\_stress}
  sont implémentées mais non appelées ailleurs dans les fichiers transmis.
  \item \textbf{L'étape d'agrégation du pipeline n'est pas câblée} entre
  \texttt{run\_scenarios.py} et \texttt{monte\_carlo.agreger}.
  \item \textbf{Écart de critère d'arrêt avec M4}, documenté en section~4.2
  et non encore arbitré par le groupe.
  \item \textbf{Attribution de \texttt{exp6\_at\_risk\_districts.py}
  ambiguë}, développée en détail en section~15.
  \item \textbf{Recouvrement possible entre l'Expérience 2 (attendue de
  M5/M6) et la Figure~3 de l'Expérience 6 (M3)}, toutes deux portant sur la
  vérification empirique de la décroissance en $1/\sqrt{N}$.
  \item \textbf{Paramètre sensible : $\mu$.} L'Expérience 3 documente
  explicitement que $\mu$ contrôle un compromis précision/vitesse ; un choix
  de $\mu$ mal justifié dans le rapport final fragiliserait la comparaison
  avec la stratégie de référence.
  \item \textbf{Hypothèse forte : coûts strictement positifs.} La convexité
  stricte de $J$ dépend de $c_e > 0$ pour tout $e$ ; aucun contrôle
  explicite de cette hypothèse n'apparaît dans \texttt{objective.py} lui-même
  (elle serait normalement garantie en amont, par M1/M2).
  \item \textbf{Chargement silencieux d'une configuration manquante.} La
  classe \texttt{DemandModel} de \texttt{demand\_model.py} retourne un
  dictionnaire vide si \texttt{network\_config.json} n'est trouvé à aucun des
  deux chemins testés, plutôt que de lever une exception explicite --- ce qui
  contraste avec le choix assumé ailleurs dans \texttt{demand\_model.py} de
  ne jamais avoir de repli silencieux.
  \item \textbf{Absence de tests transmis.} Impossible de confirmer l'état
  réel de validation de \texttt{descente\_projetee\_par\_lot} face à
  \texttt{descente\_projetee}, ni de \texttt{projeter\_orthant\_positif}.
  \item \textbf{Absence de résultats numériques transmis.} Aucune valeur
  réelle de $\kappa$, de nombre d'itérations, de violation résiduelle ou de
  coût technique n'est disponible pour cette fiche (section~11).
\end{enumerate}

% =====================================================================
\section{Interface avec le Membre 6}
% =====================================================================

Cette section synthétise, à l'intention de M6, l'état exploitable de la
contribution de M5.

\begin{description}[style=nextline]
  \item[Modules directement utilisables] \texttt{objective.py} et
  \texttt{gradient\_descent.py} sont complets, cohérents entre eux et
  exercés par trois expériences fonctionnelles (3, 4, 5). Ce sont les
  modules les plus fiables de la contribution de M5 à ce stade de l'audit.
  \item[Fonctions constituant des interfaces importantes]
  \texttt{descente\_projetee(A, b, couts, mu, \ldots)} pour résoudre un
  problème à une demande donnée ; \texttt{cout},
  \texttt{violation\_contrainte} pour évaluer une solution ;
  \texttt{pas\_maximal\_theorique(A, couts, mu)} pour obtenir la borne $2/L$
  nécessaire à tout choix de pas.
  \item[Données récupérables] Historiques de convergence complets
  (\texttt{HistoriqueConvergence}) pour les Expériences 3, 4, 5 ; résultats
  structurés par scénario (\texttt{ResultatScenario}) si M6 exerce
  \texttt{run\_scenarios.py}, ce qui n'a pas encore été fait par les
  fichiers fournis.
  \item[Résultats déjà disponibles] Aucune valeur numérique concrète n'a été
  transmise avec les fichiers audités (voir section~11). M6 doit obtenir de
  M5 le contenu réel de \texttt{results/tables/*.csv} et
  \texttt{results/figures/*.png} pour intégrer des chiffres au rapport.
  \item[Expériences intégrables immédiatement] Les Expériences 3, 4 et 5,
  sous réserve d'obtenir leurs sorties réelles : le code est complet et
  cohérent, les protocoles sont documentés et justifiés.
  \item[Expériences dépendantes du travail de M6] Les Expériences 1
  (comparaison stratégie de référence / solution optimale) et 2 (robustesse
  Monte-Carlo), selon les informations transmises initialement par M5. Le
  code fourni suggère que ces deux expériences mobiliseraient
  \texttt{run\_scenarios.py} et \texttt{monte\_carlo.py}, actuellement
  inutilisés par aucun script fourni.
  \item[Statut de l'Expérience 6] À traiter avec prudence : le script est
  attribué à M3 par son propre en-tête, et non à M5. M6 doit clarifier avec
  M3 et M5 lequel des deux volets (le script fourni, ou une version
  utilisant \texttt{monte\_carlo.py} et \texttt{quartiers\_atypiques}) doit
  entrer dans le rapport final, en particulier pour éviter une redondance
  avec l'Expérience 2.
  \item[Fichiers de résultats à consulter] \texttt{results/tables/
  exp3\_sensibilite\_mu.csv}, \texttt{results/tables/exp4\_convergence.csv},
  \texttt{results/tables/exp5\_maillage\_conditionnement.csv} ; aucun de ces
  trois fichiers n'a été transmis à ce stade de l'audit.
  \item[Figures intégrables au rapport] \texttt{exp3\_sensibilite\_mu.png},
  \texttt{exp4\_convergence.png}, \texttt{exp5\_maillage\_conditionnement.png}
  --- sous réserve de vérifier leur contenu réel, non transmis.
  \item[Informations méthodologiques devant apparaître dans le rapport
  final] Le choix du pas $\eta = 1/L$ par défaut et sa justification ; le
  choix du critère d'arrêt par déplacement (et l'écart assumé avec le
  document de M4, à faire trancher avant rédaction finale) ; le protocole
  retenu pour chaque expérience concernant le recalcul ou non de $\eta$
  d'une condition à l'autre (section~10) ; la distinction entre coût
  technique et terme de pénalisation dans toute comparaison avec la
  stratégie de référence (\texttt{objective.cout} est la métrique
  pertinente, pas \texttt{objective.objectif}).
\end{description}

% =====================================================================
\section{État d'avancement de M5}
% =====================================================================

\begin{longtable}{L{4.3cm}L{2.6cm}L{3.3cm}L{3.3cm}}
\caption{État d'avancement de la contribution de M5.}
\label{tab:avancement}\\
\toprule
Élément & État & Fichier & Résultat \\
\midrule
\endhead
Génération de scénarios (\texttt{generer\_scenarios}) & Terminé (implémenté) & \texttt{monte\_carlo.py} & Non exploité par les scripts fournis \\
Agrégation statistique (\texttt{agreger}, \texttt{erreur\_estimation}, \texttt{convergence\_par\_taille}) & Terminé (implémenté) & \texttt{monte\_carlo.py} & Non exploité empiriquement dans les fichiers fournis \\
Fonction objectif et gradient & Terminé, avec contrôle numérique & \texttt{objective.py} & Utilisé par exp3, exp4, exp5 \\
Descente de gradient projetée & Terminé & \texttt{gradient\_descent.py} & Utilisé par exp3, exp4, exp5 \\
Résolution par lot (\texttt{resoudre\_tous}, \texttt{descente\_projetee\_par\_lot}) & Terminé (implémenté) & \texttt{run\_scenarios.py} & Non exploité par les scripts fournis \\
Scénarios de stress (\texttt{scenarios\_de\_stress}) & Terminé (implémenté) & \texttt{run\_scenarios.py} & Non exploité par les scripts fournis \\
Expérience 3 (sensibilité $\mu$) & Terminé, table et figure produites par le code & \texttt{exp3\_sensitivity\_mu.py} & Valeurs numériques : à valider (non transmises) \\
Expérience 4 (vérification convergence) & Terminé, table et figure produites par le code & \texttt{exp4\_convergence\_check.py} & Valeurs numériques : à valider (non transmises) \\
Expérience 5 (maillage vs conditionnement) & Terminé, table et figure produites par le code & \texttt{exp5\_conditioning\_vs\_mesh.py} & Valeurs numériques : à valider (non transmises) \\
Expérience 6 (quartiers à risque) & Terminé, 4 figures produites --- attribution à confirmer avec le groupe & \texttt{exp6\_at\_risk\_districts.py} & Attribué à M3 dans le code ; à valider \\
Expérience 1 (baseline vs optimal) & Dépend d'un autre membre & \texttt{exp1\_baseline\_vs\_optimal.py} & Information non disponible \\
Expérience 2 (robustesse Monte-Carlo) & Dépend d'un autre membre & \texttt{exp2\_monte\_carlo\_robustness.py} & Information non disponible ; recouvrement possible avec la Figure~3 de exp6 \\
Tests unitaires de M5 & À valider & référencés (\texttt{test\_projection.py} et un test de comparaison boucle/lot), non transmis & Information non disponible \\
\bottomrule
\end{longtable}

% =====================================================================
\section{Checklist de transmission à M6}
% =====================================================================

\begin{itemize}[label=$\square$]
  \item Contenu réel de \texttt{results/tables/exp3\_sensibilite\_mu.csv}
  transmis.
  \item Contenu réel de \texttt{results/tables/exp4\_convergence.csv}
  transmis.
  \item Contenu réel de \texttt{results/tables/
  exp5\_maillage\_conditionnement.csv} transmis.
  \item Images des figures \texttt{exp3\_sensibilite\_mu.png},
  \texttt{exp4\_convergence.png}, \texttt{exp5\_maillage\_conditionnement.png}
  transmises.
  \item Statut de l'Expérience 6 clarifié avec M3 et le groupe (script
  conservé tel quel, réattribué, ou remplacé par une version utilisant
  \texttt{monte\_carlo.py}).
  \item Recouvrement entre l'Expérience 2 et la Figure~3 de l'Expérience 6
  tranché.
  \item Écart de critère d'arrêt avec le document de M4 (section~4.2)
  arbitré et harmonisé entre le code et le rapport.
  \item Fichiers de tests réels (\texttt{tests/test\_projection.py} et le
  test de comparaison boucle/lot) transmis ou leur absence confirmée.
  \item Contenu de \texttt{network\_config.json} transmis, pour documenter
  précisément le format d'entrée du pipeline.
  \item Contenu de \texttt{src/graph/build\_graph.py} et
  \texttt{src/graph/graph\_analysis.py} (M1/M2) obtenu si M6 a besoin de
  documenter l'interface amont de \texttt{run\_scenarios.py} et de
  l'Expérience 5.
  \item Choix méthodologique tranché sur la lecture retenue pour
  \texttt{scenarios\_de\_stress} (quantile marginal par quartier, ou
  quantile de la demande totale).
  \item Décision du groupe sur le nombre de scénarios Monte-Carlo à retenir
  pour \texttt{generer\_scenarios}, actuellement non fixé dans le code.
\end{itemize}

\end{document}
